\documentclass[11pt]{article}

\usepackage[margin=1in]{geometry}
\usepackage{hyperref}
\usepackage{graphicx}
\usepackage{amsmath, amssymb}
\usepackage{enumitem}
\usepackage{url}
\usepackage[T1]{fontenc}

\title{PALM-mbQTL: Quick User Guide}
\author{Peter Li}
\date{\today}

\begin{document}
\maketitle

\section{Overview}

PALM-mbQTL implements a Poisson mixed-effect model for microbiome QTL
analysis. For each microbial feature (e.g.\ genus abundance), the
pipeline

\begin{enumerate}[nosep]
  \item builds a genetic relationship matrix (GRM) from SNP genotypes;
  \item merges microbiome abundance with sample-level covariates;
  \item fits a null Poisson GLMM with random genetic effect; and
  \item performs single-variant score tests across the genome.
\end{enumerate}

This document describes the minimal Docker-based installation and the
example workflow in \texttt{example/}, which is driven by the shell
script \texttt{example/quick.sh}.  The same interface can be used on
real data by editing the parameter block at the top of
\texttt{quick.sh}.

An additional PALM mode based on our own implementation is also provided; a brief usage example is given below, while full details are documented in a separate PALM manual (PDF).


\section{Installation (Docker only)}

Before running PALM-mbQTL, please make sure Docker is installed and that
you have permission to run \texttt{docker} from the command line.
Docker can be installed by following the official instructions for your
operating system (Docker Desktop for Windows/macOS, or the Docker
Engine packages for Linux).

\subsection{Pull the image}

Pull the pre-built PALM-mbQTL image from Docker Hub:

\begin{verbatim}
docker pull --platform=linux/amd64 cheeselee/palmmbqtl:latest
\end{verbatim}

(On standard x86\_64 Linux you may omit the \texttt{--platform} option,
but it is safe to keep it.)

\subsection{Run the example}

Unpack the provided example folder (containing the toy dataset and
scripts) to any location on your machine, then run the example workflow with:

\begin{verbatim}
cd /path/to/example
docker run --rm --platform=linux/amd64 \
  -v "$PWD":/work \
  cheeselee/palmmbqtl:latest \
  bash -lc "cd /work && bash quick.sh"
\end{verbatim}

This command mounts the current directory into the container at
\texttt{/work}, executes \texttt{quick.sh} inside the container, and
writes all outputs into the local \texttt{output/} subdirectory of the
example folder.


\section{Example input files}

The example data live in \texttt{example/input/}:

\begin{itemize}[nosep]
  \item \texttt{abd.txt}: microbial abundance (phenotype) table.
  \item \texttt{cov.txt}: sample-level covariates.
  \item \texttt{geno.bed}, \texttt{geno.bim}, \texttt{geno.fam}:
        PLINK binary genotype files.
  \item \texttt{grm.rds} (optional): precomputed GRM in RDS format.
\end{itemize}

\subsection{\texttt{abd.txt}: microbiome abundance}

A minimal example:

\begin{verbatim}
IID	g_1	g_2	g_3
DKR276	3789	674	728
DLB658	3404	1670	194
NTY571	1399	1269	184
\end{verbatim}

Requirements:

\begin{itemize}[nosep]
  \item Tab-delimited text file (here the extension is `.txt`).
  \item One row per sample.
  \item First column: unique sample identifier (name specified by
  	exttt{sampleIDColinabdFile}; in the example it is `IID`).
  \item Remaining columns: non-negative integer counts for microbial
    features.
  \item Column names (e.g.\ \texttt{g\_1}) are used in
    	exttt{phenoCol}.
\end{itemize}

\subsection{\texttt{cov.txt}: covariates}

Example:

\begin{verbatim}
IID	AGE	SEX	PC01	PC02	PC03	PC04	PC05
DKR276	1.82742057825961	1	0.010171	-0.00258488	-0.00222028	-0.00296875	0.00288836
DLB658	-0.305812631038716	0	-0.00853818	-0.00585994	-0.000248601	-0.0078163	0.00122123
NTY571	1.24733084590656	1	-0.00688391	-0.00522521	0.00244565	-0.00342203	0.00128414
\end{verbatim}

Requirements:

\begin{itemize}[nosep]
  \item Tab-delimited.
  \item One row per sample; sample ID column must match
    the `IID` column in `abd.txt` (specified by
  	exttt{sampleIDColincovFile}).
  \item This example does not include a sequencing depth / offset
    column. If you wish to use a log-offset (e.g. sequencing
    depth), add a column (for example `SeqDepth`) and set
  	exttt{offsetCol} accordingly; the pipeline will use
    \(\log(\texttt{offsetCol})\) as the Poisson offset.
  \item Columns listed in \texttt{covarColList} (for example
    `AGE`, `SEX`, `PC01`, `PC02`, ...) will be included as fixed
    effects in the model.
\end{itemize}

\subsection{Genotypes: \texttt{geno.bed/.bim/.fam}}

Standard PLINK 1 binary format:

\begin{itemize}[nosep]
  \item \texttt{geno.bed}: genotype matrix.
  \item \texttt{geno.bim}: SNP annotations (chromosome, SNP ID,
        position, alleles).
  \item \texttt{geno.fam}: sample information; IDs must match the
        microbiome/covariate files.
\end{itemize}

In the example, there are five samples and three SNPs.

\subsection{Example outputs}

The example runs produce two sets of outputs under \texttt{example/output/}:

\begin{itemize}[nosep]
  \item Files prefixed with \texttt{palm1} correspond to the PALM run, for
    example \texttt{palm1\_step1\_allpheno.rda} and
  	\texttt{palm1\_step2\_chr1\_g\_1.txt}.
  \item Files prefixed with \texttt{palm2} correspond to the
    PALM-mbQTL run, for example \texttt{palm2\_step1\_allpheno.rda}
    and \texttt{palm2\_step2\_chr1\_g\_1.txt}.
\end{itemize}

Example output snippets:

\paragraph{PALM (\texttt{palm1})}
\begin{verbatim}
SNP	est	stderr
chr1.1119172.G.A	-0.0654061415202243	0.0654061379291642
chr1.1423343.G.T	0.190386805592012	0.190386809336004
chr1.3522693.A.G	-0.11205843777915	0.112058432554074
\end{verbatim}

Columns: \texttt{SNP} (variant identifier), \texttt{test} (estimated
effect for the tested covariate) and \texttt{stderr} (its standard
error).

\paragraph{PALM-mbQTL (\texttt{palm2})}
\begin{verbatim}
CHR	SNP	cM	POS	A1	A2	N	AF	SCORE	VAR	PVAL
1	chr1:1119172:G:A	0	1119172	A	G	100	0.91	-1.24601	4.16174	0.541347
1	chr1:1423343:G:T	0	1423343	T	G	100	0.865	-3.73529	5.45663	0.10981
1	chr1:3522693:A:G	0	3522693	G	A	100	0.925	1.97883	5.00944	0.376629
\end{verbatim}

Columns: \texttt{CHR}, \texttt{SNP} (variant id), \texttt{POS} (bp
position), \texttt{A1}/\texttt{A2} (alleles), \texttt{N} (sample
  size), \texttt{AF} (allele frequency for \texttt{A1}), \texttt{SCORE}
  (score statistic), \texttt{VAR} (variance of the score), and
  	\texttt{PVAL} (single-variant p-value).

\section{\texttt{quick.sh}: parameter block}

The script \texttt{example/quick.sh} is the main user interface.  At
the top of the file there is a small block where all user-facing
parameters are defined:

\subsection*{Required parameters}

\begin{description}[style=unboxed,leftmargin=0cm]
  \item[\texttt{PALMmethod}] Which workflow to run.
    \begin{itemize}[nosep]
      \item \texttt{1}: run the \texttt{PALM} workflow (\texttt{step1\_palm.R} + \texttt{step2\_palm.R});
      \item \texttt{2}: run the \texttt{PALM-mbQTL} workflow (\texttt{step0\_generateGRM.R} + \texttt{step1\_fitNULL.R} + \texttt{step2\_scoreTest.R}).
    \end{itemize}

  \item[\texttt{inputFolder}] Path (relative to \texttt{quick.sh})
    containing all input files. In the example it is \texttt{input}.

  \item[\texttt{outputFolder}] Directory for all outputs; will be
    created if it does not exist.

  \item[\texttt{genoFile}]
    Genotype input. If \texttt{genoFile} is given without a file extension (e.g.\ \texttt{./input/geno}),
    the script treats it as a PLINK prefix and expects the binary files
    \texttt{.bed}, \texttt{.bim}, and \texttt{.fam}. 
    % If \texttt{genoFile} ends with \texttt{.gds}, it is interpreted as a single GDS genotype file.
    Now only PLINK format is supported.
    
  \item[\texttt{abdFile}] Path to the microbiome abundance table (example: \texttt{abd.txt}).

  \item[\texttt{covFile}] Path to the covariate file (example: \texttt{cov.txt}).

  \item[\texttt{sampleIDColinabdFile}] Column name of the sample ID in \texttt{abdFile} (example: \texttt{IID}).

  \item[\texttt{sampleIDColincovFile}] Column name of the sample ID in \texttt{covFile} (example: \texttt{IID}).

  \item[\texttt{chrom}] Chromosome to analyze in step~2. This value is passed to the step~2 script via \texttt{--chrom=\$\{chrom\}}.

  \item[\texttt{covarColList}] Covariates to be included as fixed effects (used in \texttt{PALMmethod=2} step~1 via \texttt{--covarColList}).
    In the example it is a comma-separated list \texttt{AGE,SEX}.

\end{description}

\subsection*{Optional parameters}

\begin{description}[style=unboxed,leftmargin=0cm]
  \item[\texttt{grmFile}] Path to the GRM RDS file (used only when \texttt{PALMmethod=2}).
    If the file does not exist, step~0 generates it from \texttt{genoFile};
    otherwise the existing GRM will be reused.

  \item[\texttt{palm1\_step1\_prefix}] Output prefix for step~1 when \texttt{PALMmethod=1}.
    The null model is saved as \texttt{\$\{palm1\_step1\_prefix\}.rda}.

  \item[\texttt{palm2\_step1\_prefix}] Output prefix for step~1 when \texttt{PALMmethod=2}.
    The null model is saved as \texttt{\$\{palm2\_step1\_prefix\}.rda}.

  \item[\texttt{palm2\_step2\_prefix}] Output path/prefix for step~2 results when \texttt{PALMmethod=2}.
    The default uses Bash parameter expansion \texttt{\$\{chrom:+\_chr\$\{chrom\}\}}
    so that \texttt{chrom=1} produces a suffix like \texttt{\_chr1}.
\end{description}



\section{Workflow steps}

Internally, \texttt{quick.sh} orchestrates four core R scripts that are
bundled inside the Docker image (under \texttt{extdata/} in the source
repository):

\begin{enumerate}
  \item \textbf{Generate GRM from genotypes}
  (\texttt{step0\_generateGRM.R})\\
  Takes the PLINK/GDS files defined by \texttt{genoFile} and produces a
  GRM saved to \texttt{grmFile}.  This step only needs to be run once
  per genotype dataset.

  \item \textbf{Fit null GLMM}
  (\texttt{step1\_fitNULL.R})\\
  For each phenotype in the list \texttt{PHENOS}, fits a Poisson
  mixed-effect model with
  \begin{itemize}[nosep]
    \item log-offset \(\log(\texttt{offsetCol})\),
    \item fixed covariates from \texttt{covarColList}, and
    \item random effect with covariance proportional to the GRM.
  \end{itemize}
  The fitted model object is saved as
  \texttt{\$outputFolder/\$pheno\_step1.rda}.

  \item \textbf{Score-based association test}
  (\texttt{step2\_scoreTest.R})\\
  Given the null model and genotype data, performs single-variant
  score tests across all SNPs in \texttt{genoFile}.  Results are
  written to
  \texttt{\$outputFolder/\$pheno\_step2.txt}.

   \item \textbf{Generate feature information summary}
  (\texttt{step3\_info.R})\\
  Summarizes the microbiome abundance table \texttt{abdFile} and writes a
  feature-level report. The report includes
  \begin{itemize}[nosep]
    \item \texttt{FeatureID}: all feature/phenotype columns in \texttt{abdFile} (excluding the sample ID column),
    \item \texttt{Prevalence}: fraction of samples with non-missing values for each feature, and
    \item \texttt{AvgProportion}: mean abundance/proportion (NA removed).
  \end{itemize}
  In addition, the script reads the PLINK family file
  \texttt{\$genoFile.fam} to append dataset-level metadata to the end of the report:
  total number of samples and total number of families (unique \texttt{FID}).
  This step is typically run once per abundance/genotype dataset for QC and reporting.
  
\end{enumerate}

As a user, you typically do not need to call these R scripts
individually; instead, you:

\begin{enumerate}[nosep]
  \item edit the parameter block at the top of
        \texttt{example/quick.sh};
  \item run \texttt{bash quick.sh} inside the Docker container; and
  \item inspect the per-phenotype outputs under \texttt{output/}.
\end{enumerate}


Besides the Poisson mixed-model mbQTL pipeline, PALM-mbQTL also provides
a built-in \texttt{PALM} mode, implemented in this project, for testing
associations between microbial features (e.g.\ genus-level abundances)
and sample-level covariates without using genotype data. This mode is
driven by \texttt{example/quick\_palm.sh}, which calls the internal R
function \texttt{palm()} and outputs a table of feature–covariate
association results. A separate document (PALM PDF) describes the PALM
model, input format, and all available arguments in detail.


\section{Notes and extensions}

\begin{itemize}[nosep]
  \item The example uses only a few samples and SNPs to keep the
        runtime short. For real analyses, simply replace the
        example input files with your full dataset and adjust the
        parameter block accordingly.
  \item This document focuses on quick-start shell wrappers
        (\texttt{quick.sh}, \texttt{quick\_palm.sh})
        for running the pipeline end-to-end. If you would like to
        deploy PALM-mbQTL step by step on an HPC scheduler (e.g. SLURM
        or HTCondor), or integrate it into existing container
        workflows (Docker/Apptainer/Singularity \texttt{.sif} images),
        please contact the author for assistance.
\end{itemize}


\end{document}
